\documentclass[paper=a4, fontsize=11pt]{article}

%----------------------------------------------------------------------------------------
%	PACKAGES AND OTHER DOCUMENT CONFIGURATIONS
%----------------------------------------------------------------------------------------

\usepackage{amsmath,amsfonts,amsthm} % Math packages
\usepackage{sectsty} % Allows customizing section commands
\allsectionsfont{\centering \normalfont\scshape} % Make all sections centered, the default font and small caps

\usepackage{fancyhdr} % Custom headers and footers
\pagestyle{fancyplain} % Makes all pages in the document conform to the custom headers and footers
\fancyhead{} % No page header - if you want one, create it in the same way as the footers below
\fancyfoot[L]{} % Empty left footer
\fancyfoot[C]{} % Empty center footer
\fancyfoot[R]{\thepage} % Page numbering for right footer
\renewcommand{\headrulewidth}{0pt} % Remove header underlines
\renewcommand{\footrulewidth}{0pt} % Remove footer underlines
\setlength{\headheight}{13.6pt} % Customize the height of the header

\numberwithin{equation}{section} % Number equations within sections (i.e. 1.1, 1.2, 2.1, 2.2 instead of 1, 2, 3, 4)
\numberwithin{figure}{section} % Number figures within sections (i.e. 1.1, 1.2, 2.1, 2.2 instead of 1, 2, 3, 4)
\numberwithin{table}{section} % Number tables within sections (i.e. 1.1, 1.2, 2.1, 2.2 instead of 1, 2, 3, 4)

%\setlength\parindent{0pt} % Removes all indentation from paragraphs - comment this line for an assignment with lots of text
\setlength{\parskip}{1em}
\renewcommand{\baselinestretch}{1.5}

\usepackage{tikz-qtree}
\usepackage{lscape}

\usepackage{graphicx}
\graphicspath{ {images/} }

\usepackage{xepersian}
\settextfont[Path=fonts/]{Vazir.ttf}
%\setlatintextfont{Times New Roman}

%----------------------------------------------------------------------------------------
%	TITLE SECTION
%----------------------------------------------------------------------------------------

\newcommand{\horrule}[1]{\rule{\linewidth}{#1}} % Create horizontal rule command with 1 argument of height

\title{
\normalfont\normalsize
\includegraphics[scale=0.1]{aut}
\hspace{5cm}
\includegraphics[scale=0.1]{ceit} \\
\textsc دانشگاه صنعتی امیرکبیر \\
\textsc دانشکده مهندسی کامپیوتر و فناوری اطلاعات
\horrule{0.5pt} \\ [0.4cm] % Thin top horizontal rule
\huge ارزیابی کارآیی سیستم‌های و شبکه‌های کامپیوتری \\ % The assignment title
\huge تمرین دوم \\ % The assignment title
\horrule{2pt} \\ [0.5cm] % Thick bottom horizontal rule
}

\author{پرهام الوانی}

\date{\normalsize\today} % Today's date or a custom date

\begin{document}

\maketitle % Print the title

\section{سوال اول}

\begin{align}
\begin{split}
    f^{(1)}_{11} = p^{(1)}_{11} = a \\
    f^{(2)}_{11} = p^{(2)}_{11} - f^{(1)}_{11}p^{(1)}_{11} = a^2 - a * a = 0 \\
    \cdots\\
    f_{11}^{(n)} = 0 \quad \forall n \ge 2
\end{split}
\end{align}

بنابراین
\lr{state}
اول \lr{transient} است. 

\begin{align}
\begin{split}
    f^{(1)}_{22} = p^{(1)}_{22} = 0 \\
    f^{(2)}_{22} = p^{(2)}_{22} - f^{(1)}_{22}p^{(1)}_{22} = p^{(2)}_{22} = 1 * 1 = 1 \\
    f^{(3)}_{22} = p^{(3)}_{22} - f^{(1)}_{22}p^{(2)}_{22} - f^{(2)}_{22}p^{(1)}_{22} = p^{(3)}_{22} = 0 \\
    \cdots \\
    f_{22}^{(n)} = 0 \quad \forall n \ge 3 \\
    \sum_{n=1}^{\infty} f_{22}^{(n)} = f_{22}^{(2)} = 1 \\
    \sum_{n=1}^{\infty} nf_{22}^{(n)} = 2 * f_{22}^{(2)} = 2 < \infty \\
\end{split}
\end{align}

بنابراین
\lr{state}
دوم \lr{positive recurrent} است.

با همین روند \lr{state} سوم نیز
\lr{positive recurrent} است.

با توجه به این موضوع که:

\begin{align}
\begin{split}
    f_{21}^{(n)} = f_{31}^{(n)} = 0 \quad n \ge 1
\end{split}
\end{align}

از وضعیت‌های ۲ و ۳ نمی‌توان به وضعیت ۱ رسید.

\section{سوال دوم}

از آنجایی که:

\begin{align}
\begin{split}
    \frac{1}{p + q}
    \begin{bmatrix}
        q & p \\
        q & p
    \end{bmatrix}
    \begin{bmatrix}
        1 - p & p \\
        q & 1 - q
    \end{bmatrix}
    =
    \frac{1}{p + q}
    \begin{bmatrix}
        q & p \\
        q & p
    \end{bmatrix}
\end{split}
\end{align}

\begin{align}
\begin{split}
    \lim_{n\to\infty} P^n = &\\
    \lim_{n\to\infty}
    & \begin{bmatrix}
        1 - p & p \\
        q & 1 - q
    \end{bmatrix}^n
    = \frac{1}{p + q}
    \begin{bmatrix}
        q & p \\
        q & p
    \end{bmatrix}
\end{split}
\end{align}

و در نهایت توزیع حدی برابر است با:

\begin{align}
\begin{split}
    \begin{bmatrix}
        q/(p + q) & p/(p + q) \\
        q/(p + q) & p/(p + q)
    \end{bmatrix}
    \begin{bmatrix}
        \alpha & 1 - \alpha
    \end{bmatrix}
    =
    \begin{bmatrix}
        q/(p + q) & p/(p + q)
    \end{bmatrix}
\end{split}
\end{align}

از آنجایی که توزیع حدی از وضعیت اولیه مستقل است، این مقادیر حالت پایه نیز می‌باشند.

\section{سوال سوم}
\section{سوال چهارم}


\end{document}